\xiaosi

\subsection{结论}

本文围绕"可解释自动驾驶规划"这一主题，提出了基于神经电路策略（NCP）的轻量可解释规划方法，并通过进化搜索自动优化其布线拓扑。主要贡献和结论如下：

（1）发现并解决了NCP在强化学习Q值估计中的表达能力瓶颈。通过将NCP全部神经元隐状态经两层MLP映射为Q值（Q-head解耦方案），在保持液态ODE动力学可解释性的同时，将性能从接近随机策略提升至与最优基线持平。

（2）NCP以7,152个参数在高速公路场景上取得了34.94的平均奖励，与拥有13,189参数的MLP（34.88）相当，参数效率是GRU的3倍、LSTM的5倍。与同参数量的全连接LTC相比，NCP的结构化稀疏拓扑使奖励提升23\%，定量证明了生物启发布线结构的优势。

（3）设计了面向NCP布线拓扑的安全感知进化搜索算法，发现最优拓扑仅需15个神经元（较手工设计减少40\%），且命令层无需循环连接。搜索结果表明NCP的ODE动力学本身已具备足够的短时记忆能力。

（4）NCP的四层结构化设计赋予了每层神经元明确的功能角色，使决策过程可被逐层追踪和审计，为自动驾驶系统的安全合规提供了结构性支撑。

\subsection{展望}

本文的工作存在以下不足和改进方向：

（1）搜索仅在单一环境上评估适应度，导致搜索出的拓扑在其他场景上泛化不足。未来可采用多环境联合搜索或元学习策略改善跨场景泛化能力。

（2）当前实验平台为轻量级的Highway-env仿真器，场景复杂度有限。未来可在CARLA或nuPlan等高保真仿真器上验证，或在真实驾驶数据上进行评估。

（3）可解释性分析目前以定性观察为主，未来可引入神经元功能归因、关键时刻激活对比等定量指标，建立更严格的可解释性评价体系。

（4）NCP的碰撞率仍高于MLP，可进一步探索安全奖励塑形、约束强化学习、或将NCP与安全过滤器结合的方案来改善安全性。

（5）ODE展开数消融实验（4.3.3节）揭示了默认超参数并非最优，但受限于时间仅在单一种子上完成，未来应在多种子上系统调优\texttt{ode\_unfolds}参数，并探索其与训练稳定性、梯度传播之间的深层关系。初步结果表明3步展开是一个值得深入研究的候选配置。

（6）可探索将闭式连续时间网络（CfC）替代LTC作为NCP的基础单元，消除ODE求解器的计算开销，同时保留液态动力学的自适应特性。
